\section{Problem statement and Goals}

The dataset consists of 56 distinct German fire-brigade floor plans carrying
1633 polygon annotations over 19 functional room categories, partitioned by a
frozen project-aware 35/10/11 train/validation/test split in which no building
project crosses a partition boundary. The class frequencies are strongly
long-tailed. Categories 1--10 (B\"uroraum, WC, Treppe, Lager, Technikraum,
Aufzug, Andere, K\"uche, ``Flur, Halle'', Besprechungsraum) hold 1551 of the 1633
annotations. The six categories designated primary-rare --- 11 Allg.\
Verkehrsfl\"ache, 12 Abstellraum, 13 Versammlungsraum, 14 Schacht, 15 Garderobe,
16 Sanit\"arraum --- hold 70 annotations in total, of which 30 fall in the
training split and 17 in the test split. Two further categories, 17 Werkstatt
and 18 Speiseraum, are labelled low-support and reported separately; category 19
Labor has zero test instances and is excluded from every quantitative claim.
The per-split counts are given in Section~\ref{sec:dataset} and plotted in
Figure~\ref{fig:classdist}.

A detector trained on this distribution learns the head classes and assigns
essentially no probability mass to the tail, so the room types of greatest
operational interest --- shafts, storage rooms, assembly rooms --- are located but
never named, or not detected at all. The affected parties are the archive owners
who would use such a model to convert legacy plans into structured building
data, and any downstream fire-safety or facility-management query that depends on
room type rather than room geometry.

Since the drawings are line art on white ground, the interior of a single room
can be repainted by a diffusion inpainting model while every wall stroke, door
symbol, text label and neighbouring annotation is preserved exactly. This makes
the dataset an unusually clean testbed for the hypothesis that generative
augmentation can fill a long tail. The project was formulated to test that
hypothesis under conditions strict enough that a negative answer would be
informative: generation must read training-split pixels only, so that no
synthetic sample can leak validation or test information; every arm must differ
from the baseline in exactly one controlled respect; and the reported effect must
survive repetition over seeds and over evaluation splits. The problem is
therefore not ``build the best floor-plan segmenter'' but ``measure, with a
defensible error bar, whether inpainting augmentation helps, and whether that
answer is a property of the augmentation or of the backbone''.

Two evaluation views are used throughout. The \emph{multiclass} view is the
19-category instance-segmentation problem scored with COCO mask AP
\cite{lin2014coco}, plus a rare-macro AP averaged over categories 11--16 with
non-zero test support. The \emph{class-agnostic} view maps every annotation to a
single \texttt{room} category and measures whether the model finds room
boundaries at all, independent of naming. Separating the two is what allows a
change in localisation to be distinguished from a change in classification, and
it turned out to be the decisive measurement of the project.

\subsection{Goals}

\subsubsection{Goal 1: Leakage-safe augmented datasets and a controlled experiment matrix}

Build a generation stack that produces additional rare-room training examples
without ever reading a validation or test image, and wrap it in an experiment
matrix in which every arm is comparable to the unaugmented baseline. Success is
verifiable rather than statistical: every generator must be provably restricted
to the frozen training image identifiers; every accepted sample must pass
automated geometry, pixel-quality and exact-structural-invariance checks with a
recorded decision audit; every arm must use the identical model, optimiser
schedule and optimizer-step budget, and must be evaluated on the untouched test
annotations. The matrix must isolate the treatment of interest --- diffusion
inpainting --- against non-generative controls, so that a measured effect can be
attributed to the generated pixels rather than to the surrounding pipeline.

\subsubsection{Goal 2: Measure the augmentation effect on Mask R-CNN}

Quantify the effect of each augmentation arm on a torchvision Mask R-CNN, in
both the multiclass and class-agnostic views, with an error bar large enough to
support or reject a claim. This requires two independent repetitions: a grid over
three random seeds $\{42, 123, 2026\}$ for all fourteen arms, and a rotated
held-out cross-validation that re-evaluates the frozen checkpoints on five
disjoint folds of the 21 held-out images, raising the pooled rare test support
from 17 to 40 instances. Success is a statistically defensible AP delta relative
to the unaugmented arm e0 --- explicitly including the outcome that no
defensible delta exists, in which case the reported conclusion is a null result
with its $n$, standard deviation and $p$-value stated.

\subsubsection{Goal 3: Test whether a foundation-model backbone changes the verdict}

Determine whether the Mask R-CNN verdict is a property of the augmentation or of
the model class, by repeating the comparison on Meta's SAM 3 with its official
trainer, open-vocabulary text-concept head and the unchanged split. The three
most informative augmentation datasets from Goal 2 are re-used verbatim, so the
only difference between the two backbones is the backbone. Success is a direct
statement of whether the sign and magnitude of the augmentation effect are
preserved, plus a measurement of the backbone effect itself, which serves as the
scale against which every augmentation delta is judged.

\subsubsection{Goal 4: Test a two-stage DINOv2 pipeline and inpainting crops as classifier augmentation}

Separate localisation from naming by pairing a class-agnostic detector with a
frozen DINOv2-base classifier over crop features, and test whether the
inpainting-pipeline rare crops help the classifier stage where they did not help
the detector. Two evaluations are required: a ground-truth-region top-1
classification that removes detector error entirely and therefore isolates the
classification bottleneck, and an end-to-end two-stage mask AP that can be
compared directly against the fine-tuned detectors on the same test split.
Success is a diagnosis --- whether the residual error is a localisation problem
or a classification problem, and whether synthetic crops move the classification
number in either direction.
